\section{讲解笔记：【第13次作业笔记】 矩阵位移法(1)}
\begin{analysisbox}
本节按题目、答案和讲解笔记顺序整理。对原图中的结构几何关系，正文保留 可辨识文字，并在图形页中保留清理后的题图，便于核对杆件、支座、荷载和尺寸。
\end{analysisbox}
\subsection*{第 1 页转写}
Fe- keze

1.基本概念

A.杆端位移和杆端力

B.杆端位移和结点位移

C.杆端位移和结点荷载

D.结点位移和结点荷载

8-2判断题。结构刚度矩阵反映了结构结点位移与荷载之间的关系。（$\times$）（烟台

结点,荷载

大学2013)

2.刚度矩阵的计算

试用先处理法建立图8－43所示结构的结构刚度矩阵，

8-3

忽略杆件的轴向变形

（EI=常数）。（湖南大学2008）

8-4

求图8－44所示刚架结构刚度矩阵的任意两个主元素（自由结点）。已知各杆

I/A=1/10。（西南交通大学2010）

L|A= I / /o

4(00.))

[0,90

( 2)

M,0

6m

4m

12EI

6EI

12EI

6ED

T12E7

6EI

2EI

6EI

4EI

ye

6EI

12EI

12EI

6E1

13

2EI

6E1

6E1

4EI

4EN

【4EI

2E1

2EI

4EN

2EI

2E1

4EI

2E1

4EI

2E1

4EI

1112

\subsection*{第 2 页转写}
h-8

T|A= 1 / /0

EA.

A= lo1

M,e

139

4m

1566

kob= 2E1

T|A= 1 / /0

M,

4m

3.等效结点荷载

8-5

图8－45所示结构按矩阵位移法计算，则与结点位移1、2（正方向见图虚线标

示）对应的等效结点荷载向量为

：（浙江大学2012）

8-6图8-45所示结构，已知两杆EI=常数。若按矩阵位移法求得结点位移山2分

别为61、$\triangle$2，则杆件BC左、右端的弯矩各为

（浙江大学2012）

8-7

试求图8-46所示结构中结点3的综合结点荷载列阵F3=（X3

（中南大学2009）

\subsection*{第 3 页转写}
4m

1/2

1/2

4m

4m

图8-45

图8-46

50

12EI

12EI

6E1

6E

199

6EI

4EI

6E1

2EI

12EI

12EI

6EI

6EI

19

13

13

199

6EI

4EI

6EI

2E1

注意：在矩阵位移法中，

如果题目中完全没有刚度相关的数据，

这种类型的题目，

是需要考虑轴向变形

的。

12kN/m

15kN-m

20kN

(w. 5.6)

(7.8.5)

(1.2.3)

p? =Pe+ Po3

60

1(0p)

4m

4m

4m

10

16

\subsection*{第 4 页转写}
8-8

按先处理法求图8－47所示连续梁的结点荷载列阵P。（西南交通大学2007）

8-9

图8-48（a）所示结构，不考虑

轴向变形，整体坐标见图8－48（b），图中

20kN-m

(1,0,2)

('01)

括号内数码为结点总码（力和位移按水平、

5kN

3kN

竖直、转动方向顺序排列），求等效结点荷载

列阵P。（山东建筑大学2008）

6kN/m

2kN

4kN

5kN-m

12kN/m

M,0

(0,0,0)

M,0

(0,0,0)

2EI

4m

4m

4m

4m

(a)

(b)

图8-47

图8-48

16

-5 m

20kN-m

(1,0,2)

(1,0,3)

5KN

3kN

注意：等效节点荷载向量，就是杆端力向量反号。

6kN/m

12

M,

(0,0,0)

(0,0,0)

4m

(a)

(b)

\&=-sd= D, Sin$\alpha$=

\subsection*{第 5 页转写}
PO-TI P :

12

-12

:-

12

8-10

图8－49所示平面结构用矩阵位移法计算，引入支承条件后的整体刚度矩阵为

多少阶？求结点2和结点5的综合结点荷载列阵。（中南大学2011）

4.内力和位移的计算

[4516)

(7.8.19)

8-11已知图8-50所示连续梁的结点

位移列阵4=（0.0030.0040.005)T（分

别为结点2、3、4的转角），设各杆的E=2$\times$

2(10)

(on)

10*kN/cm ，I=400cm，1=400cm。求单元

(2)的剪力和弯矩。（华南理工大学2010）

-q/8

4kN-m

2.4kN/m

5.6kN-m

1 [010,0)

1/2

1/2

图8-50

图8-49

92

一元

110

Q= 90$^\circ$, 0d: 0, Sind:I

\subsection*{第 6 页转写}
8-11

已知图8-50所示连续梁的结点

位移列阵$\triangle$=（-0.0030.0040.005)T（分

别为结点2、3、4的转角），设各杆的E=2$\times$

10*kN/cm ，I=400cm*，l=400cm。求单元

El= 2x10*kaltn'x ych'$\times$ 10*

4kN-m

2.4kN/m

5.6kN-m

08

EI-

kn.m

- m

IZET

12$\times$ 8w

300

.x

150

= 150

- b. 003

3.2

6$\times$8w

300

4.8

nt-

4+ 810

150

-310

32

0.004

3.2

mt

4.8

4.8

2.4

kn.m

(Q2)= 5 1 v

-1.2

\subsection*{第 7 页转写}
8 - 12

（上海大学2008、武汉大学2008类似）

2(1)

FP

1010

Fo = o I

k.b=P

4m

(0,0)

图8-51

$\alpha$= 90$^\circ$,

jd= U, (in$\alpha$=l

Ko= TI fo T-

(ind.

,

@- TT @T:

\subsection*{第 8 页转写}
22/68

576 7

5:76 7

22168

176

176

Fo = O =

-2.12

12

(0. 651>

12

9 37.04

37.00

\subsection*{第 9 页转写}
-707

0.6113

8-13

图8-52所示结构，已知结点2的水平、竖向和转角位移分别为0、一0:

一0.1667，结点3的转角为0.2220，试求23杆2端杆端力。已知单元刚度矩阵中的部芬

\subsection*{第 10 页转写}
12XEx

元素：

13

1kN

1kN-m

lo=2m

I1m4

M,e

1m

4m

图8-51

图8-52

12EI

6E1

4EI

k32=k62= ks3=k65=

k33=2k63=k66=

k22= k52=k55

（浙江大

73

学2008)

7 - 0. 0926

39

39-

- 01667

-6E

IZE

-6

自助

-66

0111.0

0. 00 1E

8 - 14

用矩阵位移法计算并作图8-53所示连续梁的弯矩图。各杆EI相同，为常数。

（重庆大学2009）

(00, 0)

(013,3)

10kN/m

30kN-m

(0,0,0)

20kN

10kN

[0,0,p)

YB

6m

3m

3m

-3

230

at

12

13

注意：如果只要求求解杆端弯矩，单元刚度矩阵可以写成2阶的。

31

